\section{Resultados}

Na nova execução, AMALIA obteve médias superiores às de Llama nas cinco métricas (Tabela~\ref{tab:resultados-globais}). O juiz classificou como corretas 67,3\% das respostas de AMALIA e 60,0\% das de Llama; EM foi 13,0\% e 6,2\%. A diferença entre EM e as métricas semânticas pode refletir, em parte, respostas válidas com formulações diferentes. O BERTScore usa Albertina; estes resultados não devem ser confundidos com os ficheiros históricos diretamente em \path{results/}.

\begin{table}[ht]
\centering\small
\begin{tabular}{lrrrrr}
\toprule
Modelo & EM & Contains & ROUGE-L & BERTScore & Juiz \\
\midrule
AMALIA & 0,130 & 0,411 & 0,334 & 0,696 & 0,673 \\
Llama 3.1 8B & 0,062 & 0,360 & 0,270 & 0,665 & 0,600 \\
\bottomrule
\end{tabular}
\caption{Médias globais da execução \texttt{principal-ptpt}, com 453 entradas por modelo e BERTScore calculado com Albertina.}
\label{tab:resultados-globais}
\end{table}

Por origem, o juiz deu a AMALIA 64,3\% em HUM e 70,4\% em LLM; a Llama, 58,7\% e 61,4\%. São diferenças observadas, não evidência de causalidade. Por tema, AMALIA atingiu 100\% em Motorizações e Energia Automóvel (20 perguntas), mas 25\% em Tradições Académicas de Coimbra (20). Os mínimos de Llama foram 30\% em Basquetebol e Tradições Académicas de Coimbra (20 cada). Com apenas 18–20 perguntas por tema, estas médias são indicativas.
